\documentclass[14pt,a4paper]{extarticle}

% =========================================================
% КОМПИЛЯТОР: XeLaTeX
% VS Code: убедитесь, что выбран рецепт XeLaTeX
% =========================================================

% -------------------------
% Язык и шрифты
% -------------------------
\usepackage{fontspec}
\usepackage{polyglossia}
\setdefaultlanguage{russian}

% Times-подобный шрифт с кириллицей.
% Сначала пробуем Times New Roman, затем Liberation Serif,
% затем FreeSerif.
\IfFontExistsTF{Times New Roman}
    {\setmainfont{Times New Roman}}
    {
        \IfFontExistsTF{Liberation Serif}
            {\setmainfont{Liberation Serif}}
            {\setmainfont{FreeSerif}}
    }

\IfFontExistsTF{DejaVu Sans Mono}
    {\setmonofont{DejaVu Sans Mono}}
    {\setmonofont{FreeMono}}

% -------------------------
% Страница
% -------------------------
\usepackage[
    left=30mm,
    right=15mm,
    top=20mm,
    bottom=20mm
]{geometry}

\usepackage{setspace}
\onehalfspacing

\setlength{\parindent}{1.25cm}
\setlength{\parskip}{0pt}
\setlength{\emergencystretch}{2em}

\usepackage{microtype}
\usepackage{ragged2e}
\usepackage{enumitem}

% -------------------------
% Графика и таблицы
% -------------------------
\usepackage{graphicx}
\usepackage{float}
\usepackage{caption}
\usepackage{array}
\usepackage{tabularx}
\usepackage{longtable}

\graphicspath{{images/}}

\captionsetup{
    font=normalsize,
    labelsep=endash,
    justification=centering
}

% -------------------------
% Заголовки
% -------------------------
\usepackage{titlesec}

\titleformat{\section}
    {\bfseries\fontsize{14pt}{17pt}\selectfont}
    {\thesection}
    {0.8em}
    {}

\titleformat{\subsection}
    {\bfseries\fontsize{14pt}{17pt}\selectfont}
    {\thesubsection}
    {0.8em}
    {}

\titleformat{\subsubsection}
    {\bfseries\fontsize{14pt}{17pt}\selectfont}
    {\thesubsubsection}
    {0.8em}
    {}

\titlespacing*{\section}
    {0pt}{12pt}{12pt}

\titlespacing*{\subsection}
    {0pt}{12pt}{6pt}

\titlespacing*{\subsubsection}
    {0pt}{12pt}{6pt}

% -------------------------
% Содержание
% -------------------------
\usepackage{tocloft}

\renewcommand{\contentsname}{СОДЕРЖАНИЕ}
\setlength{\cftbeforesecskip}{0pt}
\setlength{\cftbeforesubsecskip}{0pt}

% -------------------------
% Списки
% -------------------------
\setlist[itemize]{
    label=--,
    leftmargin=1.25cm,
    itemsep=0pt,
    topsep=3pt
}

\setlist[enumerate]{
    leftmargin=1.25cm,
    itemsep=0pt,
    topsep=3pt
}

% -------------------------
% Ссылки
% -------------------------
\usepackage[hidelinks]{hyperref}

% =========================================================
% ДАННЫЕ ДЛЯ ТИТУЛЬНОГО ЛИСТА
% =========================================================

\newcommand{\StudentGroup}{421702}
\newcommand{\StudentName}{А.\,В.~Сайковский}
\newcommand{\TeacherName}{Е.\,Д.~Кветко}
\newcommand{\VariantNumber}{19}

\newcommand{\DisciplineName}{Объектное моделирование интеллектуальных систем}
\newcommand{\LabNumber}{2}
\newcommand{\LabTitle}{Реализация объектной модели}
\newcommand{\SystemTitle}{Электронная система учета и управления личным бюджетом}

% =========================================================
% ВСПОМОГАТЕЛЬНЫЕ КОМАНДЫ
% =========================================================

% Заглушка для диаграммы.
% Пока изображения нет, можно оставить эту команду.

\newcommand{\DiagramPlaceholder}[1]{
    \begin{figure}[H]
        \centering
        \fbox{
            \parbox[c][6cm][c]{0.85\textwidth}{
                \centering
                #1\\[0.5cm]
                \textit{}
            }
        }
        \caption{#1}
    \end{figure}
}

% Пример вставки готовой диаграммы:
%
% \begin{figure}[H]
%     \centering
%     \includegraphics[width=0.95\textwidth]{er-diagram.png}
%     \caption{ER-диаграмма программной системы}
%     \label{fig:er}
% \end{figure}

% =========================================================
% ДОКУМЕНТ
% =========================================================

\begin{document}

\justifying
\pagestyle{plain}

% =========================================================
% ТИТУЛЬНЫЙ ЛИСТ
% =========================================================

\begin{titlepage}
\thispagestyle{empty}

\begin{center}

Министерство образования Республики Беларусь

Учреждение образования

\vspace{0.2cm}

\textbf{БЕЛОРУССКИЙ ГОСУДАРСТВЕННЫЙ УНИВЕРСИТЕТ}\\
\textbf{ИНФОРМАТИКИ И РАДИОЭЛЕКТРОНИКИ}

\vspace{0.3cm}

Факультет информационных технологий и управления

Кафедра интеллектуальных информационных технологий

\vfill

{\Large\textbf{ОТЧЁТ}}

\vspace{0.4cm}

о выполнении лабораторной работы №\LabNumber

по дисциплине «\DisciplineName»

\vspace{0.4cm}

на тему:

\vspace{0.3cm}

{\large\textbf{\LabTitle}}

\vspace{0.4cm}

\textbf{Объектная модель программной системы}\\
«\SystemTitle»

\vspace{0.3cm}

(Вариант №\VariantNumber)

\vfill

\end{center}

\begin{flushright}
\begin{minipage}{0.48\textwidth}

Студент гр. \StudentGroup

\StudentName

\vspace{0.6cm}

Преподаватель

\TeacherName

\end{minipage}
\end{flushright}

\vfill

\begin{center}
Минск 2026
\end{center}

\end{titlepage}

\setcounter{page}{1}

\section*{Постановка задачи}

\begin{enumerate}
    \item Найти ошибки и недочёты в спецификации системы учёта личного бюджета. Для проверки использовать список из приложения~3.
    \item Реализовать объектную модель и веб-интерфейс с переходами между экранами. Разместить исходный код на GitHub и указать ссылку в отчёте.
\end{enumerate}

\section*{Реализация задачи}

\subsection*{Задание 1. Поиск ошибок}

\begin{enumerate}
    \item \textbf{Категория операции.} В текстовом описании среди полей операции необязательным назван только комментарий. На диаграмме у операции может быть от нуля до одной категории. В спецификации не сказано явно, можно ли сохранить операцию без категории и куда попадут такие расходы в отчёте.

    \item \textbf{История лимитов.} На диаграмме у категории может быть не больше двух лимитов. В тексте ограничение относится к действующим лимитам: по одному на каждую валюту. Если пользователь закроет лимит и позже создаст новый, в истории появится ещё одна запись. Модель не отделяет действующие лимиты от архивных и не уточняет, что происходит с историей после удаления категории.

    \item \textbf{Состояния лимита.} В тексте предупреждение появляется при достижении 80\,\%. На диаграмме указано поле «состояние», но не перечислены его значения. Также не описано, как отображается лимит, когда расходы составляют ровно 100\,\%, и как рассчитывается сумма перерасхода.

    \item \textbf{Итоги за месяц.} По сценарию пользователь видит итоги отдельно в BYN и USD. Метод на диаграмме получает месяц, но не получает валюту и возвращает один набор показателей с одной валютой. В таком виде из сигнатуры неясно, как получить итоги для обеих валют. У результатов анализа по категориям и общей статистики указан тип \texttt{Map}, но состав возвращаемых данных не раскрыт.
\end{enumerate}

\subsection*{Задание 2. Реализация проекта}

Saldo --- веб-приложение для учёта личного бюджета. Серверная часть написана на Python с FastAPI, данные хранятся в PostgreSQL, а интерфейс написан на React и TypeScript. После входа пользователь может перейти к операциям, категориям, лимитам, накоплениям и аналитике. Администратору доступны блокировка и разблокировка учётных записей и общая статистика по системе.

\subsubsection*{Детали реализации}

\begin{itemize}
    \item \textbf{Вход и выход.} \texttt{AuthController} регистрирует пользователя, проверяет пароль и выдаёт JWT. При выходе браузер удаляет токен. Отдельная серверная сессия в приложении не хранится.

    \item \textbf{Операции и категории.} \texttt{Operation} хранит доходы и расходы. Операцию можно сохранить без категории. При удалении категории связанные операции сохраняются без неё, а действующие лимиты архивируются. В спецификации удаление разрешено только для категорий без связанных операций.

    \item \textbf{Лимиты.} \texttt{BudgetController} рассчитывает исполнение по запросу. Для закрытых периодов используется сохранённый снимок. При изменении лимита сохраняются данные прошлых месяцев, а при удалении --- также текущего. Предупреждение появляется при расходах строго больше 80\,\%, а превышение фиксируется выше 100\,\%. Пользователь видит сумму расходов, лимит и процент использования.

    \item \textbf{Накопления и аналитика.} В \texttt{Goal} хранится целевая и уже отложенная сумма. Пользователь обновляет накопления вручную и видит их на шкале прогресса. Отдельной отметки о достижении цели нет. \texttt{AnalyticsController} считает доходы, расходы и финансовый результат по выбранной валюте, сравнивает месяцы и группирует расходы по категориям. Операции без категории попадают в отдельную группу.

    \item \textbf{Связь с диаграммой классов.} Операциями, категориями, целями и учётными записями управляют обработчики FastAPI. Доступ к данным разделён между репозиториями, а экраны представлены компонентами React.
\end{itemize}

Исходный код проекта доступен в \href{https://github.com/Qudef1/budget_tracker}{репозитории Saldo на GitHub}.

\end{document}
